\section{The Duty of Care and the Procedural Regulation}

This chapter will provide a critical analysis of the impact of the Procedural Regulation on the respect of duty of care of \glspl{dpa} in the integrated enforcement system of the \gls{gdpr}.
As laid out in \hyperref[sec:DoC]{section 2.3.1}, the duty of care requires administrations to be careful and diligent in the way they handle complaints and investigations. 
The administration must ensure it is careful in establishing the facts and points of law relevant to a decision, so as to not produce harm to the parties involved.
This chapter will first examine issues with the pre-reform cooperation mechanism through real-world examples and academic literature.
Then it will explore how the Procedural Regulation may address the identified shortcomings, and whether it introduces new challenges or unintended consequences.

\subsection{The Deficiencies of GDPR Enforcement}
This section will provide an overview of the deficiencies in the enforcement of the \gls{gdpr} that have been found by the academic literature as well as real world cases.

\subsubsection{Case Study 1: Twitter Ireland}

\paragraph{Background}
The case concerns a bug on Twitter Inc.'s platform that allowed users to view the private tweets of other users.\autocite[7]{edpbTwitter}
That bug was reported by Twitter International Company -- subsidiary of Twitter Inc. incorporated in Ireland -- to the Irish \gls{dpa} (hereinafter '\gls{dpc}') through a data breach notification.
The breach, according to Twitter, affected more than 88.000 users in the \gls{eu}.\autocite[7]{edpbTwitter}
The Irish \gls{dpc} started an ex-officio investigation, concluding, in its draft decision to the \glspl{sac} that Twitter had violated Article 33(1) and (5) \gls{gdpr}, concerning the notification procedure by a controller of a data breach.\autocite[7]{edpbTwitter}

\paragraph{Objections by the SACs}
The \glspl{sac} raised objections to the draft decision, both on the findings of the \gls{dpc} relating to Article 33 \gls{gdpr}, as well as raising new issues that were not addressed in the draft relating to the fundamental principles of data processing laid down in article 5 \gls{gdpr}.\autocite[s 6.2]{edpbTwitter}
The \gls{dpc} rejected these objections stating that it `exercised its discretion' to restrict the scope of the investigation to the two points of Article 33 \gls{gdpr}.\autocite[92]{edpbTwitter}
It further substantiated that changing the scope of the draft at this stage of the procedure would have jeopardized the validity of the final measure within the realm of its national law, since the \gls{dpc} had informed Twitter of the scope of the investigation at the outset.\autocite[93]{edpbTwitter}

\paragraph{The EDPB's Decision}
While the \gls{edpb} was able to acknowledge that some of the objections raised by the \glspl{sac} were indeed reasoned and relevant, and they demonstrated a potential risk for the fundamental rights and freedoms of data subjects, it ultimately could not reach a decision on those points.\autocite[130]{edpbTwitter}
The \gls{edpb} concluded that it did not have enough elements to establish further infringements of the \gls{gdpr}, and recognized that the \gls{dpc}'s decision to restrict the scope of inquiry directly affects the Board's ability to examine the objections raised by the \glspl{sac}.\autocite[132--133]{edpbTwitter}
The \gls{edpb} reminded that an \gls{lsa} should frame an inquiry in a way that allows the \glspl{sac} to effectively fulfill their role in determining whether an infringement of the \gls{gdpr} has occurred.\autocite[135--136]{edpbTwitter}

\subsubsection{Case Study 2: WhatsApp Ireland}

\paragraph{Background}
The case concerns an ex-officio investigation by the \gls{dpc} on WhatsApp Ireland Limited, a company with its single establishment in Dublin.\autocite[1]{edpbWhatsApp}
The investigation focused on whether WhatsApp Ireland had complied with Articles 12, 13 and 14 \gls{gdpr}.
The decision of the \gls{dpc} was prompted by various complaints by data subjects, although the \gls{dpc} made it clear that it was not a procedure relating to complaints, but just an investigation out of its own volition.\autocite[2]{edpbWhatsApp}

\paragraph{Objections by the SACs}
The \glspl{sac} raised objections on the specific points of law that the \gls{dpc} had already focused on, as well as on the scope of the investigation, and the infringement of Article 5(1)(a) \gls{gdpr} laying down the principle of `lawfulness, fairness and transparency' of data processing.\autocite[159--161]{edpbWhatsApp}
The \gls{dpc} similarly argued that it was in its own remit to manage its own procedure, and that different investigations on other points of law were being conducted parallel to the one at hand.\autocite[163--164]{edpbWhatsApp}

\paragraph{The EDPB's Decision}
Ultimately, the \gls{edpb} rejected the objections raised by the \glspl{sac} on the scope of the investigation, as it found that they were not sufficiently reasoned and did not identify how and which questions and issues should have been examined.\autocite[171]{edpbWhatsApp}
However, in the case of Article 5(1)(a) \gls{gdpr}, the assessment was more complex, and reproduces some of the dynamics seen in the Twitter case.
The Italian \gls{sa} argued that a finding of infringement of Article 5(1)(a) \gls{gdpr} would derive from a violation of Articles 12 to 14 \gls{gdpr}, and that the \gls{dpc} should ultimately find the controller in breach of that principle of data processing.\autocite[184]{edpbWhatsApp}
The \gls{edpb} acknowledged that the principle of `transparency' in 5(1)(a) \gls{gdpr} is of broader nature than the obligations laid down in Articles 12 to 14 \gls{gdpr} and other substantive obligations.\autocite[192]{edpbWhatsApp}
The \gls{edpb} ultimately concluded from the draft decision itself that the \gls{dpc} had already found that the nature of the infringements committed by WhatsApp Ireland were reflective of a `significant level of non-compliance', and thus ordered the \gls{dpc} to find the controller in breach of Article 5(1)(a) \gls{gdpr}.\autocite[197--201]{edpbWhatsApp}

\subsubsection{Analysis}

Both cases demonstrate some of the challenges that the \gls{gdpr}'s cooperation and consistency mechanism has faced in practice.
In both procedures, the \glspl{sac} raised objections to the scope of the investigation carried out by the \gls{lsa}, arguing that it was too narrow and did not allow for a comprehensive assessment of the potential infringements of the \gls{gdpr}.

While in the case of WhatsApp the \gls{edpb} rejected the objections on scope for not being sufficiently reasoned, the Twitter case outright showed how a lack of information and evidence leads authorities to be unable to substantiate potential infringements.
The \gls{edpb} recognized that the reduced scope (and therefore, information gathering) of the investigation by the \gls{dpc} directly affected the possibility of concretely enforcing the law to its full extent.
As outlined above, a relevant and reasoned objection outlines, in particular, a risk for the fundamental rights and freedoms of data subjects if the draft decision were to be adopted.
The protection of that fundamental right is the very purpose of the \gls{gdpr}, as it is the practical implementation in secondary law of the right to data protection enshrined in Article 16 \gls{tfeu}.
In that sense, the \gls{dpc}'s decision to restrict the scope of the investigation can be connected with a failure to respect the duty of care, as it may have produced harm to data subjects whose privacy was invaded by not fully investigating the potential infringements of their rights.

\subsection{The Procedural Regulation's Impact on the Duty of Care}

As outlined in section 3, the Procedural Regulation introduces many new provisions that aim to address the deficiencies in the cooperation mechanism of the \gls{gdpr}.
This section will provide an analysis of how those provisions may impact the respect of the duty of care by the integrated enforcement system of the \gls{gdpr}.

\subsubsection{Early Resolution}
The early resolution of complaints seems particularly problematic from the perspective of the duty of care, and the rule of law more generally.
If a complaint concerns an infringement of the rights of a data subject as laid out in ch III \gls{gdpr}, the \gls{dpa} receiving a complaint or the \gls{lsa} may opt to resolve the complaint if the infringement is brought to an end.
While one of the stated objectives of the Procedural Regulation is to make the enforcement of the \gls{gdpr} more efficient, and while timeliness is also an aspect of the right to good administration and of the duty of care, it is questionable whether the early resolution mechanism could be abused by \glspl{dpa} to avoid fully investigating infringements to the rights of data subjects.

Early resolution must be triggered at the earliest stages of the procedure, escaping the new requirements for the presentation of preliminary findings or a summary of key issues. 
The \gls{sa} receiving the complaint may decide to declare a complaint devoid of purpose on its own, before even involving the \gls{lsa}.
The \gls{lsa} may also decide to resolve the complaint if the infringement is brought to an end before the presentation of preliminary findings in the complex procedure, or in the simple procedure before the presentation of a draft decision.
When a complaint is resolved by the \gls{lsa}, a draft decision must be prepared by it.
However, the Procedural Regulation prescribes that relevant and reasoned objections by the \glspl{sac} cannot concern the early resolution of a complaint, disallowing the \glspl{sac} to raise objections on the use of the early resolution mechanism.

Granted, the early resolution procedure requires the \gls{sa} to substantiate the end of the infringement with evidence, and an \gls{lsa} retains the right to start proceedings if it considers it appropriate.
However, it becomes incumbent solely on the complainant to challenge the decision to resolve the complaint and substantiate why the complaint is not `devoid of purpose'.
It is unclear why a complainant, when requesting the exercise of their rights conferred to them by the \gls{gdpr}, should be burdened with the task of challenging a decision by a \gls{dpa}, when they have notified the authority of a violation of the law.
Specifically, \glspl{sa} are allowed to require complainants to first request the exercise of their rights from the controller as a condition for submitting a complaint with the authority.
This presupposes that a data subject has already made a request to the controller, that controller has not complied with the request, thus violating the law, and that the data subject has taken it upon themselves to report the violation to the \gls{sa}.

There is a risk that the early resolution mechanism could be used by controllers to evade their obligations under ch III \gls{gdpr}, by bringing their infringement to an end only when they are faced with the potential for an investigation by a \gls{dpa}.

\subsubsection{Simple Procedure}
The simple procedure, as outlined in section 3.4.3 largely resembles the pre-reform cooperation mechanism, excluding most of the new procedural steps and time limits from being applicable to it.
Compared to early resolution, the simple procedure is only used where there could be no reasonable doubt on the scope of the investigation, and where additional cooperation between the \gls{lsa} and the \glspl{sac} is not necessary particularly when the case can be addressed based on past similar decisions.
The simple procedure can also be discarded at the request of the \glspl{sac} with seemingly no proof.

It is important to note that under the simple procedure, the \glspl{sa} do not have to go through the new steps on consensus finding through the presentation of a summary of key issues, nor do they have the possibility of referring the initial divergences on consensus to the \gls{edpb} for a binding decision.
New procedural rights and opportunities for the parties under investigation and complainants to be heard are also not applicable to the simple procedure, with the Regulation opting instead for leaving it to the discretion of the national procedural laws of the \gls{lsa}.

As an exception to the new rules on relevant and reasoned objections, the \glspl{sac} are not prohibited from raising objections on the scope of the investigations when objecting to the draft decision.
However, as outlined in the cases above, this reproduces the same issues present with the pre-reform cooperation mechanism, where the scope of the investigation was contested at the end, and the \gls{edpb} was not able to reach a decision on the objections due to a lack of information and evidence.

\subsubsection{Early consensus finding under the complex procedure}

As noted in the cases, some of the failures to enforce the \gls{gdpr} to its full extent were due to a lack of information and evidence, and a mismatch of scope between the investigation by the \gls{lsa} and the objections raised by the \glspl{sac}.
The Procedural Regulation introduces formal procedural steps to ensure consensus on the scope of the investigation to be carried out by the \gls{lsa} and the \glspl{sac}.
The \glspl{lsa} must now present, within three months from the confirmation of its competence ($T_0$) a summary of key issues to the \glspl{sac}, outlining:
\begin{itemize}
    \item the main relevant facts;
    \item a preliminary identification of the scope of the investigation, including the legal provisions allegedly infringed;
    \item the legal and factual issues identified;
    \item an analysis of the views of the parties under investigation and complainants, if available;
    \item a preliminary identification of the corrective measures.
\end{itemize}
The \glspl{sa} will then have to find consensus on these points, and should employ mutual assistance and joint operations to do so.
Consensus should be found on the scope of the investigation, the legal and factual issues, and the corrective measures.
If consensus cannot be found, the \gls{lsa} must refer the matter to the \gls{edpb} for a binding decision. 
Interestingly, the binding decision to be requested at this stage concerns the scope of the investigation, and the Procedural Regulation does not provide for a referral on the legal and factual issues or the corrective measures at this stage.

As seen in the cases above, a lack of evidence and information can lead to a failure to enforce the law, even when it is plausible that further infringements could have been investigated.
The introduction of early-consensus finding recalls one of the \gls{edpb}'s recommendations in the abovementioned cases, where it stressed the importance of seeking consensus from all \glspl{sa} involved prior to the commencement of the investigation.

As a reminder, the duty of care requires administrations to be careful in handling complaints, and to ensure their investigations are as comprehensive as possible. 
Administrations must investigate matters going beyond what the complainant raised, as to effectively protect individuals.
In this vein, the requirement for early consensus is welcome. 
While consensus does not equate perfect enforcement, as there is nothing preventing the \glspl{sa} from all agreeing on a narrow scope of investigation, even one objection to the scope may trigger a referral to the \gls{edpb} for a binding decision.
The \gls{edpb} is, nonetheless, mostly formed by the \glspl{sa} themselves, voting on a simple majority.
